\chapter{系统路线图：从来源可行到运营工作台}

\section{先看一条完整的因果链}

这个项目不是先做一个 Agent，再回头寻找它可以使用的数据。开发顺序是一条不可跳步的依赖链：

\begin{center}
来源可行性 $\rightarrow$ 冻结快照 $\rightarrow$ 限定类目 $\rightarrow$ 证据合同 $\rightarrow$
本地证据引擎 $\rightarrow$ 检索基线 $\rightarrow$ 受约束 Agent $\rightarrow$ 运营工作台
\end{center}

每一步只解决一个新的不确定性。后一步必须继承前一步的输入身份和边界，不能因为界面更完整或模型能生成流畅文本，就改写上游数据事实。

\begin{beginnerbox}
可以把这条链理解成建造一条可审计的生产线：P0 确认原料来源和工作条件；P1 封存本次使用的原料；P2 决定只做哪个业务范围；P3 规定什么可以叫证据；P4--P5 负责稳定取证和比较方法；P6 让模型在规则内协作；P7 才把能力交给运营人员试用。
\end{beginnerbox}

\section{八个阶段的输入、能力、闸门与交接}

\small
\begin{longtable}{@{}>{\raggedright\arraybackslash}p{0.06\textwidth}>{\raggedright\arraybackslash}p{0.18\textwidth}>{\raggedright\arraybackslash}p{0.20\textwidth}>{\raggedright\arraybackslash}p{0.18\textwidth}>{\raggedright\arraybackslash}p{0.18\textwidth}@{}}
\toprule
阶段 & 输入 & 新增能力 & Gate & 交给下一阶段 \\
\midrule
\endfirsthead
\toprule
阶段 & 输入 & 新增能力 & Gate & 交给下一阶段 \\
\midrule
\endhead
P0 & 官方来源说明、字段需求、本地环境 & 有界来源探测与资源预检 & 来源、环境、容量及公开边界可行 & 经核实的来源角色与限制 \\
P1 & P0 认可的评论与元数据来源 & 可恢复下载、完整性登记和冻结快照 & 两个角色完整且 manifest 原子发布 & 按角色读取的冻结输入 \\
P2 & P1 manifest & 流式质量审计与类目决策 & 精确关联、覆盖和选择理由可审查 & Ice Makers 范围及审计身份 \\
P3 & P2 范围与真实字段 & 冻结证据口径和标签来源 & 八个 aspect、极性、强度与引用规则冻结 & 带 lineage 的证据合同 \\
P4 & P1--P3 冻结身份 & SQLite/FTS5 本地证据引擎 & 精确关联、稳定引用、原子建库 & 可确定性查询的证据服务 \\
P5 & P4 引擎与冻结案例 & 同基准检索比较 & 引用、关联、无答案和延迟门 & 是否增加检索复杂度的证据 \\
P6 & P3--P5 合同、工具与基线 & 受约束单 Agent 和 Verifier & 离线、真实模型及完成审计的有界工程门 & 结构化、安全的任务结果 \\
P7 & P6 图、工具、验证器和公开 DTO & 本地运营工作台 & 工程验证完成；真人结构化试用仍开放 & 用户验证记录与下一轮产品决策 \\
\bottomrule
\end{longtable}
\normalsize

\section{状态不是业务成功证明}

\begin{factbox}
当前仓库中，P0--P6 的阶段状态为 PASSED。这个状态只表示对应冻结输入和规定 Gate 已满足，不表示销量、利润、缺陷率、用户采纳或生产 SLA 已获证明。

P7 仍为 \statusprogress，因为目标运营用户的结构化试用尚未完成。
\end{factbox}

\begin{boundarybox}
路线图中的箭头表示工程依赖，不表示因果业务效果。尤其是 P7 尚缺目标运营用户的结构化试用，因此当前不能声称工作台节省了真实工时或改善了决策质量。
\end{boundarybox}

\section{怎样使用后面的阶段章节}

P0--P7 采用同一套阅读模板：先看“承上启下”和业务问题，再沿“开发过程”理解为什么这样实现；随后学习技术知识、复盘问题与修复、核对验收边界，最后练习如何进入下一阶段和如何在面试中准确表达。初学者应按顺序阅读，因为 P6 的 Agent 权限来自 P3--P5，而不是模型自身。
